\documentclass{discussion}

\usepackage{amsmath, amsthm, amssymb}
% \usepackage{xcolor, hyperref, cleveref}
\usepackage{mathtools}
\usepackage{derivative}
\usepackage{tabularx}
\usepackage{graphicx}
\usepackage{tasks}
\usepackage{mdframed}
\usepackage{interval}
\usepackage{commands}

\settasks{
    label-width = 13.9917pt,
    item-indent = 17.64131pt,
    after-item-skip = 90pt,
    column-sep = 20pt,
    label = (\arabic*)
}

\discussion{7}
\author{Sections B04 and B06}
\date{Fall 2024}

\begin{document}

\begin{question}
    Find the vector form or the parametric form for each of the following lines in \(\real^3\).
    \begin{tasks}(2)
        \task The line passing through \((1, 1, 1)\) and parallel to the vector \((1, 2, 3)\).
        \task The line passing through the points \((1, 2, 3)\) and \((4, 5, 6)\).
        \task The intersection of the planes \(x + y + z = 1\) and \(x - y + z = 1\).
        \task The line passing through \((1, 0, 0)\) and intersecting orthogonally with the line passing through \((0, 1, 0)\) and \((0, 0, 1)\).
        \task The line which is tangent to the curve \(\vec{r}(t) = \left( \cos 2t, \sin 2t, t \right)\) at \(t = \frac{\pi}{4}\).
        \task The line which is tangent to the curve \(\vec{r}(t) = \left( t, t^2, t^3 \right)\) at the point \((1, 1, 1)\).
    \end{tasks}
\end{question}

\vspace{80pt}

\begin{question}
    Match each pair of planes in \(\real^3\) with the appropriate relationship between them.

    \begin{center}
        \setlength{\tabcolsep}{2em}
        \begin{tabularx}{\textwidth}{lX}
            \(x + 2y - z = 1\) and \(-2 x - 4 y + 2z = -2\) & \linespread{1.4}\selectfont The planes are parallel with a positive distance of \rule{1.5cm}{0.15mm} \vspace{40pt} \\
            \(x + 2y + 3z = 1\) and \(x + 2y + 3z = -2\) & \linespread{1.4}\selectfont The planes are intersecting with an angle whose cosine is \rule{1.5cm}{0.15mm} \vspace{40pt} \\
            \(x + 2y + 3z = 1\) and \(3x + 2y + 1 z = 1\) & The planes coincide
        \end{tabularx}
    \end{center}
\end{question}

\end{document}